\section{True North}

\begin{quotex}
It goes without saying, that nobody initiates anybody else, if we understand by initiation the mystery of the Second Birth, the great sacrament, the initiation from above which is operative from above and has the value and duration of eternity. The Initiator is above...

\flright{\textsc{Valentin Tomberg}\footnote{\url{http://www.gornahoor.net/library/Meditations\_versionTwo.pdf}}}
\end{quotex}


This truth is not a ``Christian'' truth per se, but a ``great Truth'' or mystery \emph{inherent in the pagan sacraments themselves}. Had not Christianity elucidated them, doubtless some other higher vehicle or divine pattern would have impressed itself upon the heated wax of the human element which had emerged around the Mediterranean Basin in the years of Augustus, Pax Romana (for example, had Saint Paul not fulfilled his calling of being the thirteenth apostle untimely born to replace Judas). Hermeticism admits no division in the Infinite between sacred or secular (contra modern Christendom) and permits no dichotomy between natural and artificial (contra modern Paganism); rather, the Hermetic current (which is really the predestined energy of the divine) moves back and forth between earth and heaven, body and spirit, pagan and Christian, ascending the Scala Naturae\footnote{\url{https://en.wikipedia.org/wiki/Great\_chain\_of\_being}} or Jacob's Ladder, sometimes rising on a quest, sometimes descending on a mission. As Tomberg meditates, the hermit is neither active nor contemplative (nor any other such dichotomy) -- rather, he walks. He is a hermit-pilgrim.

Doubtless this is frustrating to some minds; it is, however, very illuminating to many spirits who cannot seem to find a ``home'' in the pre-arranged boxes of normal (or abnormal) life. Did Gornahoor convert me from my Christian blinkers I wore? Certainly -- however, there is a sense in which it was meant to happen regardless, and when it did, it was (as Plato teaches) a re-discovery of what was already there.

Re-reading many of Gornahoor's old posts, I was struck by how many ``answers'' have been outlined here; how quickly I forget the painstaking groundwork of the founding father(s). Is there a lesson here, about being forgetful? Does self-complacency induce automatic amnesia? Yet here, if one wills to see and accept, is enough metaphysical form and experiential datum to begin the long yet blessed trip to the True North. Let us examine a small and seemingly superficial incident in the life of the Church and England, to poetically illustrate what True North can be. We begin by reminding ourselves of the teaching, through a prism:

\begin{quotex}
To teach with kindly benevolence, not to lose one's temper and avenge the unreasonableness of others, that is the virile energy of the South that is followed by the well-bred man. To sleep on a heap of arms and untanned skins, to die unflinching and as if dying were not enough, that is the virile energy of the North that is followed by the brave man.
\end{quotex}


Evola is here quoting a Chinese sage\footnote{\url{https://en.wikipedia.org/wiki/Chung\_Yung}}, and it is worth noting that the superlative paragons of Southron and Northerner are both regarded under the rubric of excellence and true North; the Southerner could not be compassionate without the polar light. However, our genetics and earthly lives mean something.

We know that the geographical patterns of Earth mimic, mirror, and ultimately embody the imperishable realm of true spirit. Even someone as ``Christian'' as CS Lewis could grasp this concept all on his own -- but then, even Christians aren't like Lewis\footnote{\url{https://en.wikipedia.org/wiki/The\_Pilgrim\%27s\_Regress}} any more.

\begin{quotationx}
The ontological blueprint, or map of meaning, which Lewis presents is both economical and easy on the eye. Each of the cardinal directions (North, South, etc.) has a significance that is simple but allows Lewis to adequately describe the philosophical atmosphere of earth---a realm where asceticism and hedonism clash, while the supernatural exerts a haunting magnetism on either side.

``This sweet Desire cuts across our ordinary distinctions between wanting and having. To have it is, by definition, a want: to want it, we find, is to have it.''Lewis presents a vision of the world that has the memorable lines of Narnia, in that it both embodies and transcends temporal reality. The terrain of \emph{Pilgrim's Regress---}frozen North, swampy South and inexplicable ``Grand Canyon''\emph{---} will be hard to forget. All the more so because I trace the same landscapes with my spirit.\footnote{\url{https://bittersweetblue.blogspot.com/2006/07/pilgrims-regress-cs-lewis.html}}
\end{quotationx}

There is also the strange coincidence\footnote{\url{http://www.george-macdonald.com/resources/sundar\_singh.html}} of Lewis' conversion at or around the time of the Sudha's death, and Lewis' literary allusion to this fact in That Hideous Strength. Did Sudha Singh\footnote{\url{https://en.wikipedia.org/wiki/Sadhu\_Sundar\_Singh}} represent a candle lit at the time of George MacDonald's death, a candle then passed as a torch to Lewis? It is a strange thought.

It is a well-known, but poorly understood, fact of English literary and philosophical history that the group known as the Inklings (including their predecessors Chesterton and MacDonald) were engaged in a project of re-mythologizing that which had been de-mythologized, a project cut short by the War and the deaths of certain key members.

\begin{quotex}
And if we believe that God is everywhere, why should we not think Him present even in the coincidences that sometimes seem so strange? For if He be in the things that coincide, He must be in the coincidence of those things.
\end{quotex}


Although Lewis' love for things ``northern'' was certainly original enough with him, these things tend to have coincidences or connections with other events; Charles Williams' (his colleague) might have termed them ``co-inherences''. MacDonald was very interested in exploring the consciousness of man, as well as the possibility of a salvific role for imagination, which was magnetic towards true North\footnote{\url{http://www.snc.edu/english/documents/North\_Wind/By\_genre\_or\_topic/Theology/Shadows\_that\_Fall=\_The\_Immanence\_of\_Heaven\_in\_the\_Fiction\_of\_C.S.\_Lewis\_and\_George\_MacDonald\_-\_David\_Manley.pdf}}. Is it possible that something cannot be imagined to begin with (I speak here of higher imagination, rather than the modern tendency towards sexual fantasy) unless \emph{it already is}? Or does the imagination make it so, by being \emph{the already is}, and the become? The Inklings were chasing the ``holiness of the heart's affections''; what they had dimly realized was that imagination represented a living Being and a creative power.

Certainly, their philosophical background enabled this. Lewis and the other Inklings (including Williams and Barfield) were basically Platonists; they had no use for a vision of reality which excluded the golden Chain of Being\footnote{\url{https://en.wikipedia.org/wiki/Great\_chain\_of\_being}}, which used to be a uniquely Western formulation of a much older religion, and which formed the basis of well-bred men's worldviews, since Boethius and Augustine, who located the Forms in the mind of God.

To believe in a hierarchical world that was orderly, invisible, higher, and transcendent was also to accept (by necessity) that this primordial Being was immanent, effective, real, and con-substantial with Nature -- the more man came to understand this, the more he was able to embody these same ideals in a way that both diffused and intensified them. For a long, long time, the West moved in this orbit, and preserved the central truth while creating a strong civilization of altar and hearth, which persisted until the mid 19th century in dominating both the mind and the heart, at least in theory, often in practice.

It wasn't a compassion-filled, Protestant revival which gave us the only culture-group modern American Christians can call their own, but rather a neo-Platonic, Oxford educated, contemplative or hermetic, literary group of semi-celibate men, whose very brief association gave rise to enormous influence. Later on, TS Eliot would attempt to speak of the objective correlative\footnote{\url{https://en.wikipedia.org/wiki/Objective\_correlative}}, and long before, Samuel Taylor Coleridge suspected\footnote{\url{https://books.google.com/books?id=qMqmiNbjuHEC\&pg=PT493\&lpg=PT493\&dq=coleridge+analogy\&source=bl\&ots=HJOAsPwViG\&sig=9OmcyP4rqq4yewHKpzg6hQi-G-M\&hl=en\&ei=HResTumoEYq6tgexzJTgDg\&sa=X\&oi=book\_result\&ct=result\&resnum=3\&ved=0CC8Q6AEwAg\#v=onepage\&q=coleridge\%20analogy\&f=false}} many of these truths.

Lewis' \emph{Discarded Image}, Williams' \emph{The Descent of the Dove}, Tolkien's Northern myths, Barfield's \emph{Poetic Diction} -- these are emblems of the true North. If they accomplished so much with so little, what shall more focus and energy bring? Their influence grows daily, despite their death, not least where needed most -- in the dessicated Christian Church. To Christians, I say this -- if you ignore these prophets, why should God send you a greater? To pagans, I beg you to re-consider whether these men are united, by secret sympathies, with your own cause, to the glory of One at the top of the ladder?

Can we not re-unite them, though they are dead, in the memory of the heart? And were they not right to suspect that poetry, real poetry of the moral imagination, is not in conflict with true metaphysic, but with it does well consist?


\flrightit{Posted on 2011-10-29 by Logres}

\begin{center}* * *\end{center}

\begin{footnotesize}
\begin{sffamily}
\texttt{EXIT on 2011-10-29 at 13:23 said:}

``It goes without saying, that nobody initiates anybody else''

Tomberg says this as if it is some obvious truth. On the contrary, people who believe that the divine is going to descend upon them often delude themselves into thinking that some low level magical experience constitutes enlightenment. Quotes like this only attract occultist kooks and counterfeits who profane the mysteries.

\hfill

\texttt{Charlotte Cowell on 2011-10-29 at 13:50 said:}

Wonderful quote from Meditations on the Tarot, I think it's entirely true and have never forgotten it -- just one of the many profound truths contained therein. Love the image as well Cx

\hfill

\texttt{Charlotte Cowell on 2011-10-29 at 13:55 said:}

As for these, the -- `many spirits who cannot seem to find a ``home''\,' -- aren't we taught that the Son of Man has nowhere to lay His head?

Not sure how you can say Christians aren't like C S Lewis anymore -- I think you're behind with your reading, I'm certain I specifically mentioned the Inklings about a week ago. In fact I even ventured into the pub where they drink during my stint at the same university :-)

\hfill

\texttt{Charlotte Cowell on 2011-10-29 at 15:07 said:}

By the way, speaking of `original Christians', I've just learned that the sealed metal books discovered in Jordan earlier this year were verified as authentic by the Jordanian authorities over the summer:

\url{http://www.jordantimes.com/index.php?news=38498}


If you don't already know of this breathtaking find you can read one of the news stories here:

\url{http://www.telegraph.co.uk/news/religion/8423689/Could-this-couples-Bible-codices-tell-the-true-story-of-Christs-life.html}


and here

\url{http://www.dailymail.co.uk/news/article-1372741/Hidden-cave-First-portrait-Jesus-1-70-ancient-books.html}


This may be the earliest known image of Christ and these may truly BE the sealed books spoken of in Revelation, including the book that may only be opened by the Messiah.

They were found about 60 miles from Masada and about 100 from Qumran

\hfill

\texttt{Logres on 2011-10-29 at 15:58 said:}

EXIT, you are right about that danger. But existence is inherently dangerous -- you can't avoid all danger, by installing infallible safeguards. That just makes the mystery a ``secret'' -- instead, a real mystery is safe in plain sight.

Charlotte, once again, my choice of words is less than perfect; perhaps for ``home'' one could read ``a familiar place''. Thanks for your thoughts.

\hfill

\texttt{Logres on 2011-10-29 at 16:00 said:}

He doesn't, in any case, mean that there aren't people who provide initiations, just that they are not themselves the source of the Second Birth, just as he also means that the Second Birth can dispense with an initiator if ``it'' desires -- that doesn't mean that orders or preceptors aren't desirable from a human point of view.

\hfill

\texttt{HOO on 2011-10-29 at 18:31 said:}

\url{https://www.youtube.com/watch?v=QQqncOP6qd8}

``Circa Survive -- Strange Terrain''

\url{https://upload.wikimedia.org/wikipedia/en/c/cf/Blueskynoise.jpg}

\hfill

\end{sffamily}
\end{footnotesize}

